\documentclass[10pt,a4paper]{amsart}
\usepackage{amsmath,amssymb,amsthm,graphicx}
\usepackage{verbatim}
\newenvironment{reference}{\comment}{\endcomment}
%\newenvironment{reference}{}{}
\newenvironment{slogan}{\comment}{\endcomment}
\newenvironment{history}{\comment}{\endcomment}
% For commutative diagrams we use Xy-pic
\usepackage[all]{xy}

% We use 2cell for 2-commutative diagrams.
\xyoption{2cell}
% Theorem environments.
%
\theoremstyle{plain}
\newtheorem{theorem}[subsection]{Theorem}
\newtheorem{proposition}[subsection]{Proposition}
\newtheorem{lemma}[subsection]{Lemma}

\theoremstyle{definition}
\newtheorem{definition}[subsection]{Definition}
\newtheorem{example}[subsection]{Example}
\newtheorem{exercise}[subsection]{Exercise}
\newtheorem{situation}[subsection]{Situation}

\theoremstyle{remark}
\newtheorem{remark}[subsection]{Remark}
\newtheorem{remarks}[subsection]{Remarks}

\numberwithin{equation}{subsection}

% Macros
%
\def\lim{\mathop{\mathrm{lim}}\nolimits}
\def\colim{\mathop{\mathrm{colim}}\nolimits}
\def\Spec{\mathop{\mathrm{Spec}}}
\def\Hom{\mathop{\mathrm{Hom}}\nolimits}
\def\Ext{\mathop{\mathrm{Ext}}\nolimits}
\def\SheafHom{\mathop{\mathcal{H}\!\mathit{om}}\nolimits}
\def\SheafExt{\mathop{\mathcal{E}\!\mathit{xt}}\nolimits}
\def\Sch{\mathit{Sch}}
\def\Mor{\mathop{\mathrm{Mor}}\nolimits}
\def\Ob{\mathop{\mathrm{Ob}}\nolimits}
\def\Sh{\mathop{\mathit{Sh}}\nolimits}
\def\NL{\mathop{N\!L}\nolimits}
\def\CH{\mathop{\mathrm{CH}}\nolimits}
\def\proetale{{pro\text{-}\acute{e}tale}}
\def\etale{{\acute{e}tale}}
\def\QCoh{\mathit{QCoh}}
\def\Ker{\mathop{\mathrm{Ker}}}
\def\Im{\mathop{\mathrm{Im}}}
\def\Coker{\mathop{\mathrm{Coker}}}
\def\Coim{\mathop{\mathrm{Coim}}}

% Boxtimes
%
\DeclareMathSymbol{\boxtimes}{\mathbin}{AMSa}{"02}

%
% Macros for moduli stacks/spaces
%
\def\QCohstack{\mathcal{QC}\!\mathit{oh}}
\def\Cohstack{\mathcal{C}\!\mathit{oh}}
\def\Spacesstack{\mathcal{S}\!\mathit{paces}}
\def\Quotfunctor{\mathrm{Quot}}
\def\Hilbfunctor{\mathrm{Hilb}}
\def\Curvesstack{\mathcal{C}\!\mathit{urves}}
\def\Polarizedstack{\mathcal{P}\!\mathit{olarized}}
\def\Complexesstack{\mathcal{C}\!\mathit{omplexes}}
% \Pic is the operator that assigns to X its picard group, usage \Pic(X)
% \Picardstack_{X/B} denotes the Picard stack of X over B
% \Picardfunctor_{X/B} denotes the Picard functor of X over B
\def\Pic{\mathop{\mathrm{Pic}}\nolimits}
\def\Picardstack{\mathcal{P}\!\mathit{ic}}
\def\Picardfunctor{\mathrm{Pic}}
\def\Deformationcategory{\mathcal{D}\!\mathit{ef}}
\usepackage{cleveref}
\crefname{originallemma}{original lemma}{original lemmas}
\Crefname{originallemma}{Original lemma}{Original lemmas}
\makeatletter
\newlabel{cohomology-lemma-base-change-cup-product}{{0H99}{source}{}{sourcepointer-1}{}}
\newlabel{cohomology-lemma-base-change-cup-product@cref}{{[originallemma][0][]0H99}{[1][0][]source}{}{}{}}
\makeatother
\crefname{originallemma}{original lemma}{original lemmas}
\Crefname{originallemma}{Original lemma}{Original lemmas}
\makeatletter
\newlabel{duality-lemma-smooth-proper}{{0BRT}{source}{}{sourcepointer-2}{}}
\newlabel{duality-lemma-smooth-proper@cref}{{[originallemma][0][]0BRT}{[1][0][]source}{}{}{}}
\makeatother
\crefname{originalremark}{original remark}{original remarks}
\Crefname{originalremark}{Original remark}{Original remarks}
\makeatletter
\newlabel{duality-remark-relative-dualizing-complex}{{0B6S}{source}{}{sourcepointer-3}{}}
\newlabel{duality-remark-relative-dualizing-complex@cref}{{[originalremark][0][]0B6S}{[1][0][]source}{}{}{}}
\makeatother
\crefname{originalremark}{original remark}{original remarks}
\Crefname{originalremark}{Original remark}{Original remarks}
\makeatletter
\newlabel{duality-remark-relative-dualizing-complex-relative-cup-product}{{0G81}{source}{}{sourcepointer-4}{}}
\newlabel{duality-remark-relative-dualizing-complex-relative-cup-product@cref}{{[originalremark][0][]0G81}{[1][0][]source}{}{}{}}
\makeatother
\crefname{originallemma}{original lemma}{original lemmas}
\Crefname{originallemma}{Original lemma}{Original lemmas}
\makeatletter
\newlabel{lemma-proper-smooth-de-Rham}{{0FM0}{source}{}{sourcepointer-5}{}}
\newlabel{lemma-proper-smooth-de-Rham@cref}{{[originallemma][0][]0FM0}{[1][0][]source}{}{}{}}
\makeatother
\crefname{originallemma}{original lemma}{original lemmas}
\Crefname{originallemma}{Original lemma}{Original lemmas}
\makeatletter
\newlabel{lemma-relative-bottom-part-degenerates}{{0G8H}{source}{}{sourcepointer-6}{}}
\newlabel{lemma-relative-bottom-part-degenerates@cref}{{[originallemma][0][]0G8H}{[1][0][]source}{}{}{}}
\makeatother
\crefname{originallemma}{original lemma}{original lemmas}
\Crefname{originallemma}{Original lemma}{Original lemmas}
\makeatletter
\newlabel{limits-lemma-descend-dimension-d}{{0EY2}{source}{}{sourcepointer-7}{}}
\newlabel{limits-lemma-descend-dimension-d@cref}{{[originallemma][0][]0EY2}{[1][0][]source}{}{}{}}
\makeatother
\crefname{originallemma}{original lemma}{original lemmas}
\Crefname{originallemma}{Original lemma}{Original lemmas}
\makeatletter
\newlabel{limits-lemma-descend-finite-presentation}{{01ZM}{source}{}{sourcepointer-8}{}}
\newlabel{limits-lemma-descend-finite-presentation@cref}{{[originallemma][0][]01ZM}{[1][0][]source}{}{}{}}
\makeatother
\crefname{originallemma}{original lemma}{original lemmas}
\Crefname{originallemma}{Original lemma}{Original lemmas}
\makeatletter
\newlabel{limits-lemma-descend-smooth}{{0C0C}{source}{}{sourcepointer-9}{}}
\newlabel{limits-lemma-descend-smooth@cref}{{[originallemma][0][]0C0C}{[1][0][]source}{}{}{}}
\makeatother
\crefname{originallemma}{original lemma}{original lemmas}
\Crefname{originallemma}{Original lemma}{Original lemmas}
\makeatletter
\newlabel{limits-lemma-descend-surjective}{{07RR}{source}{}{sourcepointer-10}{}}
\newlabel{limits-lemma-descend-surjective@cref}{{[originallemma][0][]07RR}{[1][0][]source}{}{}{}}
\makeatother
\crefname{originallemma}{original lemma}{original lemmas}
\Crefname{originallemma}{Original lemma}{Original lemmas}
\makeatletter
\newlabel{limits-lemma-eventually-proper}{{081F}{source}{}{sourcepointer-11}{}}
\newlabel{limits-lemma-eventually-proper@cref}{{[originallemma][0][]081F}{[1][0][]source}{}{}{}}
\makeatother
\crefname{originalreference}{original reference}{original references}
\Crefname{originalreference}{Original reference}{Original references}
\makeatletter
\newlabel{situation-relative-duality}{{0G8G}{source}{}{sourcepointer-12}{}}
\newlabel{situation-relative-duality@cref}{{[originalreference][0][]0G8G}{[1][0][]source}{}{}{}}
\makeatother
\title{The relative Hodge duality pairing: existence and uniqueness of the relative trace map making the cup product pairing perfect}
\author{The Stacks Project Authors}
\date{stacks-project:derham.tex:local-snapshot-20261012; GNU FDL 1.2 or later}
\begin{document}
\maketitle
\noindent\textit{Editorial note.} All delivered mathematical text is the author's own, copied from the pinned Stacks Project chapter file derham.tex; no statement or proof step was rewritten, reordered or removed.\par\medskip
\noindent\textit{Editorial note.} Preamble provenance: the pinned chapter file opens with \textbackslash{}input\{preamble\} and carries no preamble text of its own, so the project's single Stacks-wide preamble staged at sources/stacks-derham/preamble.tex is excerpted verbatim as three spans, lines 8-12 (verbatim package and the project's own reference/slogan/history display environments), lines 14-18 (Xy-pic setup) and lines 154-217 (theorem declarations and mathematical macros). The wrapper class and the amsmath/amssymb/amsthm/graphicx base packages are supplied by the assembler. No macro was invented or redefined.\par\medskip
\noindent\textit{Editorial note.} Same-chapter pointers are located at their \textbackslash{}label lines inside the pinned chapter file; cross-chapter pointers are located in the project's own tag list staged in the snapshot (sources/stacks-chow/tags.txt), at the verbatim lines carrying their canonical tags.\par\medskip
\noindent\textit{Allowed standard prerequisites (editorial).}\begin{itemize}
\item The relative dualizing complex, its trace map, and the smooth proper duality identification with a shift of the top differentials. Duality for Schemes, remark duality-remark-relative-dualizing-complex (invoked at proof line 5899) and label duality-lemma-smooth-proper (invoked at proof line 5901). They provide the trace map \$Rf\_*\textbackslash{}omega\textasciicircum{}\textbackslash{}bullet\_\{X/S\} \textbackslash{}to \textbackslash{}mathcal\{O\}\_S\$ and, over a Noetherian base, the isomorphism \$\textbackslash{}omega\textasciicircum{}\textbackslash{}bullet\_\{X/S\} \textbackslash{}cong \textbackslash{}Omega\textasciicircum{}n\_\{X/S\}[n]\$ from which the map \$t\$ is obtained.
\item Perfectness of the relative de Rham complexes, and the relative cup product duality statement. Same chapter derham.tex, label lemma-proper-smooth-de-Rham (invoked at proof lines 5905 and 5935) and Duality for Schemes, remark duality-remark-relative-dualizing-complex-relative-cup-product (invoked at proof line 5911). The first makes \$Rf\_*\textbackslash{}Omega\textasciicircum{}p\_\{X/S\}\$ perfect, the second turns the identification \$\textbackslash{}Omega\textasciicircum{}p \textbackslash{}otimes \textbackslash{}Omega\textasciicircum{}\{n-p\} \textbackslash{}cong \textbackslash{}Omega\textasciicircum{}n\$ into the perfection of the pairing induced by the relative cup product.
\item Absolute Noetherian approximation and base change of the cup product. Limits, labels limits-lemma-descend-finite-presentation, limits-lemma-descend-smooth, limits-lemma-descend-dimension-d, limits-lemma-descend-surjective and limits-lemma-eventually-proper (invoked at proof lines 5925-5928) and Cohomology, label cohomology-lemma-base-change-cup-product (invoked at proof line 5937). They spread the situation out to a Noetherian base with a proper smooth equidimensional model, and identify the cup products of the original situation as pullbacks, so that the trace map descends to the general case.
\item The bottom part degeneration, giving a perfect complex of tor amplitude starting in degree one. Same chapter derham.tex, label lemma-relative-bottom-part-degenerates (invoked at proof line 5953). It makes \$P\$ in the distinguished triangle \$f\_*\textbackslash{}mathcal\{O\}\_X \textbackslash{}to Rf\_*\textbackslash{}mathcal\{O\}\_X \textbackslash{}to P \textbackslash{}to f\_*\textbackslash{}mathcal\{O\}\_X[1]\$ perfect of tor amplitude in \$[1, \textbackslash{}infty)\$ and the triangle locally split, which is what produces the decomposition used for uniqueness.
\end{itemize}\medskip
\ifdefined\hypertarget\hypertarget{sourcepointer-1}{}\fi
\noindent\textit{Source reference.} Original lemma 0H99: sources/stacks-chow/tags.txt, source lines 21519--21519. Navigation pointer for \textbackslash{}ref\{cohomology-lemma-base-change-cup-product\} used inside the counted span (line 5937); the argument that consumes it is retained in the span.\par\medskip
\ifdefined\hypertarget\hypertarget{sourcepointer-2}{}\fi
\noindent\textit{Source reference.} Original lemma 0BRT: sources/stacks-chow/tags.txt, source lines 14781--14781. Navigation pointer for \textbackslash{}ref\{duality-lemma-smooth-proper\} used inside the counted span (line 5901); the argument that consumes it is retained in the span.\par\medskip
\ifdefined\hypertarget\hypertarget{sourcepointer-3}{}\fi
\noindent\textit{Source reference.} Original remark 0B6S: sources/stacks-chow/tags.txt, source lines 14080--14080. Navigation pointer for \textbackslash{}ref\{duality-remark-relative-dualizing-complex\} used inside the counted span (line 5899); the argument that consumes it is retained in the span.\par\medskip
\ifdefined\hypertarget\hypertarget{sourcepointer-4}{}\fi
\noindent\textit{Source reference.} Original remark 0G81: sources/stacks-chow/tags.txt, source lines 20251--20251. Navigation pointer for \textbackslash{}ref\{duality-remark-relative-dualizing-complex-relative-cup-product\} used inside the counted span (line 5912); the argument that consumes it is retained in the span.\par\medskip
\ifdefined\hypertarget\hypertarget{sourcepointer-5}{}\fi
\noindent\textit{Source reference.} Original lemma 0FM0: sources/stacks-derham/pinned-original-source.tex, source lines 433--433. Navigation pointer for \textbackslash{}ref\{lemma-proper-smooth-de-Rham\} used inside the counted span (line 5905); the argument that consumes it is retained in the span.\par\medskip
\ifdefined\hypertarget\hypertarget{sourcepointer-6}{}\fi
\noindent\textit{Source reference.} Original lemma 0G8H: sources/stacks-derham/pinned-original-source.tex, source lines 5855--5855. Navigation pointer for \textbackslash{}ref\{lemma-relative-bottom-part-degenerates\} used inside the counted span (line 5953); the argument that consumes it is retained in the span.\par\medskip
\ifdefined\hypertarget\hypertarget{sourcepointer-7}{}\fi
\noindent\textit{Source reference.} Original lemma 0EY2: sources/stacks-chow/tags.txt, source lines 18677--18677. Navigation pointer for \textbackslash{}ref\{limits-lemma-descend-dimension-d\} used inside the counted span (line 5926); the argument that consumes it is retained in the span.\par\medskip
\ifdefined\hypertarget\hypertarget{sourcepointer-8}{}\fi
\noindent\textit{Source reference.} Original lemma 01ZM: sources/stacks-chow/tags.txt, source lines 2599--2599. Navigation pointer for \textbackslash{}ref\{limits-lemma-descend-finite-presentation\} used inside the counted span (line 5925); the argument that consumes it is retained in the span.\par\medskip
\ifdefined\hypertarget\hypertarget{sourcepointer-9}{}\fi
\noindent\textit{Source reference.} Original lemma 0C0C: sources/stacks-chow/tags.txt, source lines 15082--15082. Navigation pointer for \textbackslash{}ref\{limits-lemma-descend-smooth\} used inside the counted span (line 5926); the argument that consumes it is retained in the span.\par\medskip
\ifdefined\hypertarget\hypertarget{sourcepointer-10}{}\fi
\noindent\textit{Source reference.} Original lemma 07RR: sources/stacks-chow/tags.txt, source lines 9876--9876. Navigation pointer for \textbackslash{}ref\{limits-lemma-descend-surjective\} used inside the counted span (line 5927); the argument that consumes it is retained in the span.\par\medskip
\ifdefined\hypertarget\hypertarget{sourcepointer-11}{}\fi
\noindent\textit{Source reference.} Original lemma 081F: sources/stacks-chow/tags.txt, source lines 10215--10215. Navigation pointer for \textbackslash{}ref\{limits-lemma-eventually-proper\} used inside the counted span (line 5928); the argument that consumes it is retained in the span.\par\medskip
\ifdefined\hypertarget\hypertarget{sourcepointer-12}{}\fi
\noindent\textit{Source reference.} Original reference 0G8G: sources/stacks-derham/pinned-original-source.tex, source lines 5848--5848. Navigation pointer for \textbackslash{}ref\{situation-relative-duality\} used inside the counted span (line 5878); the argument that consumes it is retained in the span.\par\medskip
\begin{lemma}
\label{lemma-relative-duality-hodge}
In Situation \ref{situation-relative-duality} there exists an
$\mathcal{O}_S$-module map
$$
t : Rf_*\Omega^n_{X/S}[n] \longrightarrow \mathcal{O}_S
$$
unique up to precomposing by multiplication by a unit of
$H^0(X, \mathcal{O}_X)$ with the following property: for all $p$ the pairing
$$
Rf_*\Omega^p_{X/S}
\otimes_{\mathcal{O}_S}^\mathbf{L}
Rf_*\Omega^{n - p}_{X/S}[n]
\longrightarrow
\mathcal{O}_S
$$
given by the relative cup product composed with $t$
is a perfect pairing of perfect complexes on $S$.
\end{lemma}

\begin{proof}
We first assume $S$ is a Noetherian scheme.
Let $\omega^\bullet_{X/S}$ be the relative dualizing complex of $X$ over $S$ as
in Duality for Schemes, Remark \ref{duality-remark-relative-dualizing-complex}
and let $Rf_*\omega_{X/S}^\bullet \to \mathcal{O}_S$ be its trace map. By
Duality for Schemes, Lemma \ref{duality-lemma-smooth-proper}
(here we use that $S$ is Noetherian)
there exists an isomorphism $\omega^\bullet_{X/S} \cong \Omega^n_{X/S}[n]$
and using this isomorphism we obtain $t$. The complexes $Rf_*\Omega^p_{X/S}$
are perfect by Lemma \ref{lemma-proper-smooth-de-Rham}.
Since $\Omega^p_{X/S}$ is locally free and since
$\Omega^p_{X/S} \otimes_{\mathcal{O}_X} \Omega^{n - p}_{X/S} \to
\Omega^n_{X/S}$ exhibits an isomorphism $\Omega^p_{X/S} \cong
\SheafHom_{\mathcal{O}_X}(\Omega^{n - p}_{X/S}, \Omega^n_{X/S})$
we see that the pairing induced by the relative cup product is perfect by
Duality for Schemes, Remark
\ref{duality-remark-relative-dualizing-complex-relative-cup-product}.

\medskip\noindent
Proof of existence in the general case.
By absolute Noetherian approximation we can find a cartesian diagram
$$
\xymatrix{
X \ar[r] \ar[d]^f & X_0 \ar[d]^{f_0} \\
S \ar[r]^g & S_0
}
$$
with $S_0$ Noetherian and $f_0$ a proper and smooth morphism
all of whose fibres are nonempty and equidimensional of dimension $n$.
See Limits, Lemmas \ref{limits-lemma-descend-finite-presentation},
\ref{limits-lemma-descend-smooth}, \ref{limits-lemma-descend-dimension-d},
\ref{limits-lemma-descend-surjective}, and
\ref{limits-lemma-eventually-proper}. Choose a map
$$
t_0 : Rf_{0, *}\Omega^n_{X_0/S_0}[n] \longrightarrow \mathcal{O}_{S_0}
$$
as in the statement of the lemma. By cohomology and base change the
complexes $Rf_*(\Omega^p_{X/S})$ are equal to
$Lg^*(Rf_{0, *}\Omega^p_{X_0/S_0})$ for all $p \geq 0$, see
Lemma \ref{lemma-proper-smooth-de-Rham}. In particular the pullback
of $t_0$ by $g$ gives us a map $t$ as in the statement of the lemma.
Namely, by Cohomology, Lemma \ref{cohomology-lemma-base-change-cup-product},
the cup product map
$$
Rf_*\Omega^p_{X/S}
\otimes_{\mathcal{O}_S}^\mathbf{L}
Rf_*\Omega^{n - p}_{X/S}[n]
\longrightarrow
Rf_*\Omega^n_{X/S}
$$
is the pullback by $g$ of the corresponding cup product map for $f_0$.
Hence the fact that $t_0$ produces perfect pairings, pulls back to
the same property for $t$ on $S$.

\medskip\noindent
Uniqueness of $t$. Choose a distinguished triangle
$f_*\mathcal{O}_X \to Rf_*\mathcal{O}_X \to P \to f_*\mathcal{O}_X[1]$.
By Lemma \ref{lemma-relative-bottom-part-degenerates}
the object $P$ is perfect of tor amplitude in $[1, \infty)$
and the triangle is locally on $S$ split.
Thus $R\SheafHom_{\mathcal{O}_X}(P, \mathcal{O}_X)$ is perfect
of tor amplitude in $(-\infty, -1]$. Hence duality (above) shows that
locally on $S$ we have
$$
Rf_*\Omega^n_{X/S}[n] \cong
R\SheafHom_{\mathcal{O}_S}(f_*\mathcal{O}_X, \mathcal{O}_S)
\oplus R\SheafHom_{\mathcal{O}_X}(P, \mathcal{O}_X)
$$
This shows that $R^nf_*\Omega^n_{X/S}$ is finite locally free and
that we obtain a perfect $\mathcal{O}_S$-bilinear pairing
$$
f_*\mathcal{O}_X \times R^nf_*\Omega^n_{X/S} \longrightarrow \mathcal{O}_S
$$
using $t$.
This implies that any $\mathcal{O}_S$-linear map
$t' : R^nf_*\Omega^n_{X/S} \to \mathcal{O}_S$ is of the form
$t' = t \circ g$ for some
$g \in \Gamma(S, f_*\mathcal{O}_X) = \Gamma(X, \mathcal{O}_X)$.
In order for $t'$ to still determine a perfect pairing $g$ will have
to be a unit. This finishes the proof.
\end{proof}

\end{document}
